\subsubsection{Actividades [Pendientes: 1. definir casos y tolerancias de contraste, 2. concretar distribución de igniciones y escenarios, 3. fijar reglas de consulta y agregación]}\label{sec:construccion-actividades}
A partir de E3 se preparará un piloto territorial que permita examinar coherencia de entradas, comportamiento numérico y costo de ejecución, registrando tiempo, memoria, configuración y resultados. Se estudiará la sensibilidad de las cantidades de interés frente a los parámetros pertinentes y a la discretización, utilizando rangos defendibles y distinguiendo incertidumbre de entrada de efectos numéricos. Cuando existan observaciones compatibles, el contraste considerará posición del frente en tiempos comparables, dirección de propagación, área o tiempo de llegada, según lo que efectivamente permitan los datos. Un polígono final aislado no se utilizará para demostrar una trayectoria temporal completa. Con esa evidencia se concretarán las configuraciones, la transformación de clases de amenaza en espaciamientos y las combinaciones de condiciones que compondrán E4, sin extrapolar automáticamente el piloto al conjunto de los Cerros Orientales.

La aplicación se construirá integrando almacenamiento y catálogo, obtención meteorológica, búsqueda, tratamiento espacial y animación, conforme a los contratos de E2. Se comprobarán la correspondencia entre consulta y escenarios recuperados, la referencia geográfica, el orden temporal de las salidas y la exclusión del interior quemado. El procedimiento de agregación se contrastará en los casos de referencia que se acuerden para establecer qué diferencias introduce respecto de una representación defendible del incendio, documentando también situaciones desfavorables. Se registrarán por separado el desempeño de la consulta y la calidad de la aproximación, ya que una respuesta rápida no demuestra exactitud física. Las incidencias encontradas se relacionarán con la versión corregida y su nueva comprobación dentro de E6, mientras que el manual E7 se desarrollará a partir del funcionamiento real de los incrementos.
